\documentclass[a4paper, 20pt,fleqn]{article}

\usepackage{hyperref}
\usepackage{eqarrays}
\usepackage{fancyvrb}
\usepackage{graphicx}

\usepackage[utf8]{inputenc}
\usepackage{jeffe}
\usepackage{handout}
\usepackage{aurical}
\usepackage[T1]{fontenc}
\def\mle{$\le$}
\def\ith{$i^{\textrm\scriptsize th}$}

\usepackage{tikz}
\usetikzlibrary{positioning,chains,fit,shapes,calc}

\newcommand{\fillinMC}[1]{\fillinMCmath{\mbox{#1}}}
\newcommand{\fillinMCmath}[1]{\begin{tikzpicture}\draw circle [radius=0.5em];\end{tikzpicture}\ #1}
\newcommand{\fillinMCAll}[1]{\fillinMCAllmath{\mbox{#1}}}
\newcommand{\fillinMCAllmath}[1]{\begin{tikzpicture}\draw (0,0) rectangle (1em,1em);\end{tikzpicture}\ #1}
\newcommand{\fillinblank}[1]{\fillinblankmath{\mbox{#1}}}
\newcommand{\fillinblankmath}[1]{\begingroup\setlength{\fboxsep}{1em}\setlength{\fboxrule}{2pt}\fbox{\LARGE\phantom{#1}}\endgroup}
\newcommand{\answerbox}[1][3\baselineskip]{%
    \noindent\setlength{\fboxrule}{2pt}\framebox[0.4\linewidth]{%
        \raisebox{2pt}[0pt][#1]{}%
    }\par\medskip%
}
%
% The document.
%

\begin{document}
\headers{CPSC 121}{ }{Summer, 2019}

\begin{center}
    \LARGE
    \textbf{HW 5}
    \\[1ex]
    \Large Due: 23:00, Wednesday, June 19, 2019  \\
\end{center}
    \LARGE
\begin{tabular}{rl}
 & \\
CS ID 1: &  d6a2b\\
 & \\
CS ID 2: & b9i2b\\
 & \\
\end{tabular}
\large

\noindent\textbf{\LARGE Instructions:}
\begin{enumerate}
\item Do not change the problem statements we are giving you. Simply add your solutions by editing this latex document. To make it easier for the TAs to find your solutions, please use the \textbf{soln} environment we provided as follows:
\begin{soln}
\begin{verbatim}
    \begin{soln}
    My solution is here.
    \end{soln}
\end{verbatim}
\end{soln}
Your solution will then appear in blue, and be easier to differentiate from the questions.
\item Include formatting to clearly distinguish your solutions from the given problem text. Improperly or insufficiently typeset submissions will receive a penalty.
\item If you need more space, add a page between the existing pages using the \texttt{\textbackslash newpage} command.
\item Export the completed assignment as a PDF file for upload to Gradescope.
\item On Gradescope, upload only \textbf{one} copy per partnership. (Instructions for uploading to Gradescope are posted on the assignments page of the course website.)
\item During submission, for each question, please link ALL pages on which your solution appears. Submissions with several linking errors will incur a small penalty.
\item Late submissions will be accepted up to 24 hours past the deadline with a penalty of 20\% of the assignment’s maximum value
\end{enumerate}

\noindent\textbf{\LARGE Academic Conduct:} 
I certify that my assignment follows the academic conduct rules for of CPSC 121 as outlined on the course website. As part of those rules, when collaborating with anyone outside my group, (1) I and my collaborators took no record but names away, and (2) after a suitable break, my group created the assignment I am submitting without help from anyone other than the course staff. \\

\textbf{Version history:}
\begin{itemize}
    \item 2019-06-16 11:58 -- Bug corrected in Q6 ALU diagram, Q2 hint updated
    \item 2019-06-11 18:28 -- Initial version for release
\end{itemize} 

\usepackage{tikz}

\newpage
\normalsize

\newpage

\begin{Questions}
Questions 1 and 2 ask you to design a Deterministic Finite-State Automaton (DFA). You are also required to indicate the meaning of each state. Here is an example to show you what we mean by this, taken from an old assignment. The question was:

\begin{quote}
\it
  Design a  deterministic finite-state  automaton (DFA) that  accepts exactly  the strings
  over the alphabet $\{ \verb@A@, \verb@B@, \ldots, \verb@Z@ \}$ that contain at least one
  \verb@C@, at  most two $O$, and  where every $D$  comes before every $E$.  For instance,
  your DFA should accept the strings
  \begin{itemize}
  \item ACCORDEON
  \item CHOCOLATE
  \item CROCODILE
  \item SCALD
  \item SCALE
  \end{itemize}
  but not the strings
  \begin{itemize}
  \item ADRIAN (it doesn't contain the letter \verb@C@)
  \item COLONIZATION (there are three \verb@O@s)
  \item MANACLED (there is an \verb@E@ before a \verb@D@)
  \item BEWARETHEKOMODODRAGON (all of these at the same time)
  \end{itemize}
  Clearly indicate  the meaning of each  state. Hint: there are  three separate conditions
  accepted strings must meet; states will need to  encode whether or not they are met. You can label an edge with the word ``else'' to indicate it would contain every character that does not appear on another edge leaving from the same state. For instance, if a state has outgoing edges labeled ``A'', ``B'', ``E'', ``P'' and ``else'', then ``else'' stands for ``C, D, F, G, H, I, J, K, L, M, N, O, Q, R, S, T, U, V, W, X, Y, Z''.Our
  solution uses 13 states.
\end{quote}

\noindent The solution, and the explanation about the meaning of each state are listed below:
\begin{quote}
\it
  Here is a DFA that works. We used 14 states instead of 13 because having two ``garbage''
  states allows the DFA to be drawn without edges crossing. The idea is that we have three
  separate conditions to check for:
  \begin{itemize}
  \item Whether or  not a ``C'' was seen. The  states in the left half of  the picture are
    those for  which no  ``C'' has  yet been  seen. The states  in the  right half  of the
    picture are those for which a ``C'' has been seen.

  \item Whether or  not a ``E'' was  seen. The states in  the top half of  the picture are
    those for  which no  ``E'' has yet  been seen. The  states in  the bottom half  of the
    picture are those for which a ``E'' has  been seen (and hence seeing a ``D'' will lead
    the DFA to one of the garbage states).

  \item Whether we have seen no ``O'', one  ``O'' or two ``O''s. These states are drawn in
    ``layers'': the  outside rectangle  contains the  states for which  no ``O''  has been
    seen, the middle rectangle contains the states for which a single ``O'' has been seen,
    and the inner rectangle contains the states  for which two ``O''s have been seen. Once
    the DFA reaches a state where two ``O''s have been seen, any further ``O'' leads it to
    a garbage state.
  \end{itemize}
  Twelve  of  the  fourteen  states  encode  all  possible  combinations  of  these  three
  conditions.  The  remaining two  are  garbage  states that  are  used  when one  of  the
  requirements are violated (three or more ``O''s, or a ``D'' after an ``E'') and it is no
  longer possible for the DFA to accept the input string.
  \begin{center}
    \leavevmode\includegraphics[width=5in]{images/dfa-CODE.png}
  \end{center}
\end{quote}

\newpage
\end{Questions}

\begin{question}{1.}[10]

\def\StartsWith#1#2{\textit{Starts}({#1},{#2})}

\item[9] Consider the following sets and predicates: 
  \begin{itemize}
  \item $D$: video game developers
  \item $G$: video game titles
  \item $P$: hardware platforms
  \item $M(d, g)$: developer $d$ made game $g$.
  \item $R(g, p)$: game $g$ released on platform $p$.
  \item $B(d, g)$: developer $d$ bribed a journalist to write a favourable review for game $g$.
  \end{itemize}

  Rewrite each of  the following statements using \textbf{only}  the quantifiers $\forall$
  and $\exists$, the predicates $M$, $R$ and $B$, the domains $D$, $G$, $P$, and $\mathbb{R}$ (the set of real numbers),
  logical connectives, and the operators $=$, $\neq$. $<$, $\le$, $\ge$ and $>$.
  
  If you feel a  sentence is ambiguous, then state your assumed interpretation.
  \begin{question}{a.}[2]

  \item[3] Every developer made a game that released on every platform.
   
   \vspace{2.0in}
   
  \item[3] Some developer bribed a journalist to write a favourable review for a game made by a different developer.
  
  \vspace{2.0in}
   
  \item[3] Some developer bribed journalists to write favourable reviews for every game that the developer made.
  
  \vspace{2.0in}

  \end{question}

\begin{Questions}
\newpage
\end{Questions}

%
% Question 2
%
\item[9] 

Using the definitions given in question 1, translate each of the
  following predicate logic statements into English. Try to make your English translations
  as natural sounding as possible.

  \begin{question}{a.}[3]
  \item[3] $\exists d \in D, \forall x,y \in G, \exists p_1, p_2 \in P, M(d,x) \land M(d,y) \land R(x,p_1) \land R(y,p_2) \implies p_1 = p_2 \implies \lnot(\exists g \in G, M(d,g) \land B(d,g))$.
  \begin{Questions}
    \vfill
  \end{Questions}
  
  \item[3] $\exists d \in D, \forall g \in G, \lnot M(d,g) \land (R(g, \text{Nintendo Switch}) \implies B(d,g))$.
  \begin{Questions}
    \vfill
  \end{Questions}

  \item[3] $\exists d \in D, \forall p \in P, \exists  g_1, g_2 \in G, M(d, g_1) \land M(d, g_2) \land R(g_1, p) \land \lnot R(g_2, p)$.
  \begin{Questions}
    \vfill
  \end{Questions}
  
  \end{question}
\begin{Questions}
\newpage
\end{Questions}

\item[14] In this problem we will explore circuit design with limited resources by taking advantage of \textit {functionally complete} sets of logic gates. A set of logic gates is functionally complete if it is able to simulate all of the operations $\{AND, OR, NOT\}$. Here, we will show that $\{NOR\}$ alone is functionally complete.

For reference, a 2-input NOR operation on inputs $a$ and $b$ is written as $a \downarrow b$, and is equivalent to $\lnot (a \lor b)$.

\begin{enumerate}
    \item Using rules of logical equivalence, show that the 2-input NOR operation can be used to simulate a 1-input NOT operation.
    \vspace{1.0in}
    \item Using rules of logical equivalence, show that two 2-input NOR operations can be used to simulate a 2-input OR operation.
    \vspace{2.0in}
    \item Using a circuit diagram with accompanying truth table, show that multiple 2-input NOR gates (with non-inverted inputs) can be used to simulate a 2-input AND gate. Be sure to label your input and output signals.
    \vspace{2.5in}
    
    \newpage
    
    \item Below is a truth table for a system with 4 inputs. Design a circuit to implement the function, using only 2-input NOR gates with non-inverted inputs. A modest penalty will be applied to solutions using 16 or more gates. Hint: begin with a circuit using other gates, that can be easily converted to NOR gates using the double-negative law and DeMorgan's laws.
    
    \begin{tabular}{cccc|c}
      $x_3$ & $x_2$ & $x_1$ & $x_0$ & $f$ \\
      \hline
      F & F & F & F & F \\
      F & F & F & T & F \\
      F & F & T & F & T \\
      F & F & T & T & F \\
      \hline
      F & T & F & F & T \\
      F & T & F & T & F \\
      F & T & T & F & F \\
      F & T & T & T & F \\
      \hline
      T & F & F & F & F \\
      T & F & F & T & F \\
      T & F & T & F & T \\
      T & F & T & T & F \\
      \hline
      T & T & F & F & T \\
      T & T & F & T & T \\
      T & T & T & F & F \\
      T & T & T & T & F \\
    \end{tabular}
    
\end{enumerate}

\newpage

\item[8] Design  a circuit that  takes five inputs  $x_4$,  $x_3$, $x_2$, $x_1$,  $x_0$
  representing an unsigned binary number in the range of $\left[0,31\right]$ and a single
  output $z$ which is \texttt {true} if (and only if)  the binary number specified by the
  inputs is a prime number. Assume that 0 and 1 are not prime, and 2 is prime.
  For instance, your circuit should output \texttt {true} in the following cases:
  \begin{itemize}                                                                                   
  \item  $x_4 =  \texttt{false}$,  $x_3 =  \texttt{false}$, $x_2  =  \texttt{false}$, $x_1  =
    \texttt{true}$, $x_0 = \texttt{false}$ (decimal 2).
  \item  $x_4 =  \texttt{true}$,  $x_3 =  \texttt{false}$, $x_2  =  \texttt{false}$, $x_1  =
    \texttt{true}$, $x_0 = \texttt{true}$ (decimal 19).
  \item  $x_4 =  \texttt{true}$,  $x_3 =  \texttt{true}$, $x_2  =  \texttt{true}$, $x_1  =
    \texttt{false}$, $x_0 = \texttt{true}$ (decimal 29).
  \end{itemize}
  but it should output false in these cases:                                                        
  \begin{itemize}                                                                                   
  \item  $x_4 =  \texttt{false}$,  $x_3 =  \texttt{false}$, $x_2  =  \texttt{false}$, $x_1  =
    \texttt{false}$, $x_0 = \texttt{true}$ (decimal 1).
  \item  $x_4 =  \texttt{false}$,  $x_3 =  \texttt{true}$, $x_2  =  \texttt{true}$, $x_1  =
    \texttt{false}$, $x_0 = \texttt{false}$ (decimal 12).
  \item  $x_4 =  \texttt{true}$,  $x_3 =  \texttt{true}$, $x_2  =  \texttt{false}$, $x_1  =
    \texttt{false}$, $x_0 = \texttt{true}$ (decimal 25).
  \end{itemize}                                                                                     
  Justify your answer! You may use any gates available in Logisim, with any number of inputs
  and/or inverted inputs. A modest penalty will be applied to solutions using 11 or more gates.
  \newpage

\item[6] Suppose we have the following five propositions:
\begin{displaymath}
\begin{array}{l}
\phantom{\lnot}X_1 \lor X_2 \lor Y_1 \\
\lnot Y_1 \lor X_3 \lor Y_2 \\
\lnot Y_2 \lor X_4 \lor Y_3 \\
\lnot Y_3 \lor X_5 \lor Y_4 \\
\lnot Y_4 \lor X_6 \lor X_7 \\
\end{array}
\end{displaymath}

Using an \emph {indirect} proof, show that every assignment of truth values to the eleven variables $X_1$, $X_2$, $X_3$, $X_4$, $X_5$, $X_6$, $X_7$, $Y_1$, $Y_2$, $Y_3$, $Y_4$ that makes every one of the five propositions True assigns the value True to \emph {at least} one of $X_1$, $X_2$, $X_3$, $X_4$, $X_5$, $X_6$, $X_7$.
\\
\\
{\color{NavyBlue}
\\We define each statement above as a, b, c, d, e respectively.
\\ $\forall $ $X_1$, $X_2$, $X_3$, $X_4$, $X_5$, $X_6$, $X_7$, $Y_1$, $Y_2$, $Y_3$, $Y_4$ $ \in Boolean, a \land b \land c \land d \land e \rightarrow X_1 \lor X_2 \lor X_3 \lor X_4 \lor X_5 \lor X_6 \lor X_7$
\\
\\Proof:
\\ We can use proof by contradiction. Suppose $a \land b \land c \land d \land e$ is true, but $\lnot (X_1 \lor X_2 \lor X_3 \lor X_4 \lor X_5 \lor X_6 \lor X_7)$. We can use DeMorgan's theorem to change $\lnot (X_1 \lor X_2 \lor X_3 \lor X_4 \lor X_5 \lor X_6 \lor X_7)$ to $\lnot X_1 \land \lnot X_2 \land \lnot X_3 \land \lnot X_4 \land \lnot X_5 \land \lnot X_6 \land \lnot X_7$. This means that all X values has to be false. 
\\
\\Statement a is defined as $X_1 \lor X_2 \lor Y_1$, and since we know $X_1$ and $X_2$ is false, $Y_1$ has to be true for the statement to be true (property of Or gates). \\
\\Using the same logic, we can say that for true b, $Y_2$ is true; for true c, $Y_3$ is true; for true d, $Y_4$ is true. Since $Y_4$ is true, and $ X_6$ and $ X_7$ are false, statement e will evaluate to false. However, e is true for $a \land b \land c \land d \land e$ to be true. 
\\
\\Therefore, we have arrived to a contradiction, and we can conclude that the original statement of 
$\forall $ $X_1$, $X_2$, $X_3$, $X_4$, $X_5$, $X_6$, $X_7$, $Y_1$, $Y_2$, $Y_3$, $Y_4$ $ \in Boolean, a \land b \land c \land d \land e \rightarrow X_1 \lor X_2 \lor X_3 \lor X_4 \lor X_5 \lor X_6 \lor X_7$ is true.
\\
\\\boxed{}
}
\newpage

\item[6] Given a finite set $S$ of $n$ integers, $S=\{s_1, s_2, \ldots, s_n\}$, consider the problem of finding a subset of those integers, such that the sum of the items in the subset is equal to 0.

One possible solution is to generate all possible subsets of $S$, and then calculate the sum of the values in each subset to see if it is equal to zero.

\begin{question}{a.}[3]
    \item[4] The number of possible subsets for a set of $n$ integers satisfies the following equation:
    \[
          B(n)= \begin{cases}
          2 B(n-1) & \textrm{if }n > 1\\
          1 & \textrm{if }n = 0\\
          \end{cases}
    \]
    For instance, when $n=0$ there is only one possible subset: the empty set.
    When $n = 2$ there are four possible subsets:
    \[
      \begin{array}{ll}
      \{  \} &
      \{ s_1 \} \\
      \{ s_2 \} &
      \{ s_1, s_2 \} \\
      \end{array}
    \]
    Give a (very  simple) mathematical expression for $B(n)$. Use the formula from the first part to justify your answer.
    \begin{Questions}
    \vfill
    \end{Questions}
    
    \item[2] Suppose you run an algorithm which takes a number $n$ and a set $S$ of that size and then outputs every subset on a machine where your algorithm takes $.0001$ seconds for each subset it outputs. 
    
    Fill in the following table with the largest value of $n$ for which your algorithm terminates for each duration indicated.\bigskip
    
    \setlength{\extrarowheight}{0.3cm}
    \begin{tabular}{|l|c||l|c|}\hline
    Duration & Largest value of $n$ & Duration & Largest value of n \\ \hline
    1 second &  &
    1 minute & \\ \hline
    1 hour & &
    1 day & \\ \hline
    1 month & &
    1 year & \\ \hline 
    10 years & &
    100 years & \\ \hline
    \end{tabular}
    
\end{question}



\end{question}

\end{document}
